\documentclass[a4paper,11pt,reqno]{amsart}
\usepackage[T1]{fontenc}
\IfFileExists{pplr8t.tfm}{\usepackage{mathpazo}}{\usepackage{lmodern}}
\usepackage[margin=27mm]{geometry}
\IfFileExists{microtype.sty}{\usepackage{microtype}}{}
\IfFileExists{mathtools.sty}{\usepackage{mathtools}}{\usepackage{amsmath}}
\IfFileExists{booktabs.sty}{\usepackage{booktabs}}{}
\usepackage{amssymb,array}
\usepackage[hidelinks,pdfauthor={Arseniy Akopyan},pdftitle={Sixfold coincidences and obstructed convex equipartitions}]{hyperref}
\IfFileExists{enumitem.sty}{\usepackage{enumitem}\setlist[enumerate]{label=\textup{(\roman*)},leftmargin=2.3em,itemsep=3pt}}{}
\setlength{\emergencystretch}{2em}
\setlength{\parindent}{1.2em}
\allowdisplaybreaks[1]
\raggedbottom
\newcommand{\R}{\mathbb R}
\newcommand{\Z}{\mathbb Z}
\newcommand{\N}{\mathbb N}
\newcommand{\Sym}{\mathfrak S}
\newcommand{\dd}{\,d}
\newcommand{\intt}{\operatorname{int}}
\newcommand{\conv}{\operatorname{conv}}
\newcommand{\spanlin}{\operatorname{span}}
\newcommand{\codim}{\operatorname{codim}}
\newcommand{\diag}{\operatorname{Diag}}
\newcommand{\Gr}{\operatorname{Gr}}
\newcommand{\dist}{\operatorname{dist}}
\newcommand{\supp}{\operatorname{supp}}
\newcommand{\id}{\operatorname{id}}
\newcommand{\cl}{\operatorname{cl}}
\newcommand{\one}{\mathbf 1}
\newcommand{\sphere}{\mathbb S}
\newtheorem{mainthm}{Theorem}
\renewcommand{\themainthm}{\Alph{mainthm}}
\newtheorem{theorem}{Theorem}[section]
\newtheorem{lemma}[theorem]{Lemma}
\newtheorem{proposition}[theorem]{Proposition}
\newtheorem{corollary}[theorem]{Corollary}
\theoremstyle{definition}
\newtheorem{definition}[theorem]{Definition}
\theoremstyle{remark}
\newtheorem{remark}[theorem]{Remark}
\numberwithin{equation}{section}
\title[Sixfold Tverberg and convex equipartitions]{Sixfold Tverberg counterexamples
 and convex equipartition obstructions}
\author{Arseniy Akopyan}
\date{16 November 2026}
\subjclass[2020]{Primary 52C35, 55M20; Secondary 52A20, 55S91}
\keywords{Topological Tverberg theorem, convex partition, arrangement complement, equivariant map, Hausdorff-continuous functional}
\begin{document}
\begin{abstract}
We construct a counterexample for the generalized Nandakumar--Ramana Rao
problem: in every dimension $d\ge24$, there are $d-1$ continuous
functionals on convex bodies and a probability measure such that there is
no convex partition of the unit ball into six parts of equal measure on
which all these functionals are equalized.

We also show that the topological Tverberg conjecture fails for sixfold
intersections already in $\R^6$: we construct a continuous map
$\Delta^{35}\to\R^6$ whose images of six pairwise disjoint faces never
meet.
\end{abstract}
\maketitle

\section{Introduction and main results}

Let $r\ge2$ and $d\ge1$ be integers, and put $N=(r-1)(d+1)$.
Tverberg's theorem asserts that every affine map
$f:\Delta^N\to\R^d$ sends $r$ pairwise vertex-disjoint faces to
sets with a common point. The \emph{topological Tverberg conjecture}
asked whether the same holds for every continuous map. It does when
$r$ is a prime power. Frick showed that it fails for non-prime
powers \cite{Frick}, combining the higher-multiplicity Whitney trick
of Mabillard and Wagner \cite{MW} with a reduction originating in
Gromov's work \cite{Gromov}, developed as the constraint method by
Blagojevi\'c, Frick, and Ziegler \cite{BFZ}. Subsequent work gave
stronger counterexamples \cite{AMS,AKS}.

The smallest non-prime-power multiplicity is $r=6$. For a simplex
$\Delta^N$, we call a continuous map $f:\Delta^N\to\R^d$ an
\emph{almost six-embedding} if
$f(\sigma_1)\cap\cdots\cap f(\sigma_6)=\varnothing$ for every six
pairwise vertex-disjoint faces. Equivalently, no six points with
pairwise disjoint barycentric supports have the same image. Our
first result is a counterexample to the topological Tverberg
assertion in dimension six.

\begin{mainthm}\label{thm:A}
The topological Tverberg conjecture fails for sixfold intersections in
$\R^6$: there is a continuous map
\[
                         f:\Delta^{35}\longrightarrow\R^6
\]
whose images of six pairwise disjoint faces never meet.
\end{mainthm}

For $D\ge1$, let $B_D=\{x\in\R^D:\|x\|\le1\}$ be the unit ball,
let $tB_D$ be the concentric ball of radius $t>0$, and let $\mathcal K_D$
denote the Hausdorff space of nonempty compact convex subsets of $\R^D$.
A \emph{convex six-partition} of $B_D$ consists of six full-dimensional
compact convex sets covering $B_D$ with pairwise disjoint interiors.

\begin{mainthm}\label{thm:B}
In every dimension $d\ge24$, there are $d-1$ continuous functionals on
convex bodies and a probability measure on the unit ball $B_d$ such that
there is no convex partition of $B_d$ into six parts of equal measure on
which all these functionals are equalized.
\end{mainthm}

In particular, in $\R^{24}$ we obtain \emph{23 functionals and one
additional measure}. No centrality functional is concealed in this
count. Theorem~\ref{thm:B} is a consequence of a uniform obstruction
valid for every sufficiently concentrated absolutely continuous measure.

\begin{theorem}[Uniform obstruction for concentrated measures]\label{thm:concentrated}
For every integer $D\ge24$, there exist a continuous map
$\Phi_D:\mathcal K_D\to\R^{D-1}$ and a number $\eta_D\in(0,1)$ with the
following property. Let $\mu$ be any Borel measure on $\R^D$, absolutely
continuous with respect to Lebesgue measure, such that
\[
 0<\mu(B_D)<\infty,
 \qquad
 \mu(\eta_D B_D)>\frac56\mu(B_D).
\]
Then no convex six-partition $(C_1,\ldots,C_6)$ of $B_D$ satisfies both
\[
 \mu(C_i)=\frac{\mu(B_D)}6\quad(1\le i\le6),
 \qquad
 \Phi_D(C_1)=\cdots=\Phi_D(C_6).
\]
The map $\Phi_D$ and the radius $\eta_D$ are chosen before the measure
$\mu$, and are independent of it.
\end{theorem}

Let $\gamma_{D,\sigma}$ be the centered Gaussian probability measure
with covariance $\sigma^2I_D$, where $\sigma>0$.
\begin{corollary}[A smooth positive density]\label{cor:gaussian}
For every $D\ge24$, the same $\Phi_D$ from Theorem~\ref{thm:concentrated}
obstructs simultaneous equalization with $\gamma_{D,\sigma}$ whenever
$0<\sigma<\eta_D/\sqrt{6D}$. In dimension $24$, it suffices that
$0<\sigma<\eta_{24}/12$.
\end{corollary}

The fair-partition theorem of Akopyan--Avvakumov--Karasev
\cite[Theorem 7.1]{AAK} equalizes $D-1$ absolutely continuous measures
and one arbitrary continuous characteristic, for any number of parts.
Our result has a different quantifier structure: \emph{one} measure,
and $D-1$ possibly nonadditive characteristics. The two statements
are not interchangeable. Moreover, Theorem~\ref{thm:concentrated} concerns all
convex partitions, not merely normal fans or successive cuts. Its
concentration assumption is part of the statement; it does not
exclude equalization for every prescribed measure on the fixed unit
ball. A prescribed finite measure can instead be treated on a ball
whose radius is allowed to vary (Corollary~\ref{cor:prescribed}).

\subsection*{Outline of the proof}
Theorem~\ref{thm:A} is proved by scalar height tables in
Section~\ref{sec:scalar}. For the second result we prove a separate
auxiliary statement, which is the key to reducing the ambient
dimension of the convex-partition construction.

\begin{mainthm}\label{thm:C}
There is an almost six-embedding
\[
                         g:\Delta^{119}\longrightarrow\R^{22}.
\]
\end{mainthm}

Section~\ref{sec:vector} deduces Theorem~\ref{thm:C} from the
equivariant-map theorem of Avvakumov--Karasev--Skopenkov
\cite[Theorem 2.2]{AKS}, without using any higher-multiplicity
intersection-removal theorem. Joining with a constant map on a
four-simplex then gives almost six-embeddings
$\Delta^{5D-1}\to\R^{D-2}$ for $D\ge24$. In
Section~\ref{sec:cones} we show that the internal skeleta of all
six-cone partitions of $\R^D$ have parameter dimension at most
$5D-1$. Hence one can find $5D$ points that do not lie together
on such a skeleton. The resulting distance-to-boundary
functionals give Theorem~\ref{thm:concentrated} and hence Theorem~\ref{thm:B} in
Section~\ref{sec:transfer}.

Thus the historical role of the Frick--Gromov--Mabillard--Wagner
method is motivational, not an unmentioned input to our proofs.
The systematic study of spaces of convex partitions is discussed
in \cite{LZ}. Neither the target dimensions $6$ and $24$ nor the
number of functionals in Theorem~\ref{thm:B} is claimed to be
optimal.

\subsection*{Formalization status}
Both main results, Theorem~\ref{thm:A}
(\texttt{AkopyanSixfold.\allowbreak tverberg\_\allowbreak counterexample}) and
Theorem~\ref{thm:B}
(\texttt{AkopyanSixfold.\allowbreak convex\_\allowbreak partition\_\allowbreak counterexample}),
are formalized in Lean~4 (\texttt{v4.35.0}) with Mathlib~(\texttt{v4.35.0}) in
\texttt{Challenge.lean} and proved completely without \texttt{sorry} in
\texttt{Solution.lean}, using only the standard foundational axioms
\texttt{propext}, \texttt{Quot.sound}, and \texttt{Classical.choice}.
Theorem~\ref{thm:A} is verified through the explicit scalar height-table
reduction and piecewise-linear surgery in \texttt{RequestProject/};
Theorem~\ref{thm:B} is verified unconditionally for convex bodies in
$\R^d$ ($d\ge24$) under the Hausdorff metric and a Borel probability
measure (without restricting to absolutely continuous measures; Section~\ref{sec:formal}). The stronger
uniform absolutely continuous obstruction (Theorem~\ref{thm:concentrated})
and Theorem~\ref{thm:C} are not part of the checked \texttt{Challenge.lean}
targets.

\section{Height tables and their forbidden subspaces}\label{sec:tables}

Let $c,k\ge1$ be integers. Partition the vertices of $\Delta^{6c-1}$
into $c$ labeled groups, each of size six, and denote them by $v_{i,a}$,
where $i\in[c]=\{1,\ldots,c\}$ and $a\in[6]$.
For a point $x$ of this simplex, define its masses by
$\lambda_i(x)=\sum_{a=1}^6x_{i,a}$. Thus
$\lambda(x)\in\Delta^{c-1}$.

Let $\mathcal W_c=[6]^c$ be the set of words and let
$V_{c,k}=(\R^k)^{\mathcal W_c}$ be the vector space of tables.
A \emph{matching} is a six-element subset $M\subset\mathcal W_c$ whose
entries in each coordinate run through all six labels. For each matching
choose an order $m_1,\ldots,m_6$ and define a linear map
\[
 \ell_M:V_{c,k}\longrightarrow(\R^k)^5,
 \qquad
 \ell_M(z)=(z_{m_1}-z_{m_6},\ldots,z_{m_5}-z_{m_6}).
\]
Put $L_M=\ker\ell_M$ and
$Z_{c,k}=V_{c,k}\setminus\bigcup_M L_M$. The definition of $L_M$ is
independent of the chosen order. A table in $Z_{c,k}$ will be called
\emph{safe}.

A table map $H:\Delta^{c-1}\to V_{c,k}$ is \emph{compatible} if
$H(\lambda)_a$ is independent of $a_i$ whenever $\lambda_i=0$.
This is an exact boundary condition, not a limiting convention.

\begin{lemma}[Height-table realization]\label{lem:realization}
A continuous compatible map $H:\Delta^{c-1}\to Z_{c,k}$ determines an
almost six-embedding $\Delta^{6c-1}\to\R^{c+k-1}$.
\end{lemma}
\begin{proof}
For each group define a probability vector by
$p_{i,a}(x)=x_{i,a}/\lambda_i(x)$ if $\lambda_i(x)>0$, and
$p_{i,a}(x)=1/6$ otherwise. Set
\begin{equation}\label{eq:average}
 h(x)=\sum_{a\in\mathcal W_c}H(\lambda(x))_a
                         \prod_{i=1}^c p_{i,a_i}(x),
 \qquad
 f(x)=(\lambda_1(x),\ldots,\lambda_{c-1}(x),h(x)).
\end{equation}
The vector $h(x)$ belongs to $\R^k$. To check continuity at $x_0$, first
replace $H(\lambda(x))$ by $H(\lambda(x_0))$ in the finite average.
The error is at most the sup norm of their difference. In the resulting
average all factors belonging to a zero mass of $x_0$ disappear, since
the limiting table ignores their labels and each probability vector
sums to one. The remaining factors are continuous at $x_0$.

Suppose six disjoint-support points have the same $f$-image. Their
mass vectors agree. In every active group each point uses at least
one vertex; since there are only six vertices, each uses exactly one,
and all six choices are different. Complete the choices in inactive
groups by arbitrary permutations. They give a matching $M$. Compatibility
and \eqref{eq:average} identify the six heights with the six entries of
$H(\lambda)$ indexed by $M$. Their equality contradicts safety.
\end{proof}

\begin{lemma}[Separation of forbidden subspaces]\label{lem:codim}
Each $L_M$ has codimension $5k$. For distinct matchings $M,N$,
$\codim(L_M\cap L_N)\ge7k$.
\end{lemma}
\begin{proof}
The first assertion follows from the five independent vector equalities.
Two distinct matchings cannot share five words, because the missing label
in each coordinate determines the sixth word. If they share $t\ge1$
words, their joint equations identify $12-t$ entries and have codimension
$k(11-t)\ge7k$. If they are disjoint, the codimension is $10k$.
\end{proof}

We shall need one standard consequence of Alexander duality in a precise
range. All homology groups below have integer coefficients.

\begin{lemma}[Comparison by normal directions]\label{lem:comparison}
Let $L_1,\ldots,L_m$ be distinct linear subspaces of $\R^N$, each of
codimension $a\ge3$, and suppose their pairwise intersections have
codimension at least $b$, where $a+2\le b\le2a-1$. Choose surjective
normal maps $q_i:\R^N\to\R^a$ with kernel $L_i$. Then
\[
 e:\R^N\setminus\bigcup_iL_i\longrightarrow
           \prod_{i=1}^m\sphere^{a-1},
 \qquad e(z)_i=\frac{q_i(z)}{\|q_i(z)\|},
\]
is $(b-2)$-connected. Here connectivity means that the relative
homotopy groups of its mapping cylinder vanish through degree $b-2$.
For $m=1$ the map is a homotopy equivalence.
\end{lemma}
\begin{proof}
Write $Z$ for the source and $Y$ for the target. Relative general
position shows that $Z$ is $(a-2)$-connected: a filling of a sphere of
dimension at most $a-2$ can be perturbed, relative to its boundary,
to miss the finite arrangement. The same connectivity holds for $Y$.

For $m>1$, let $K=(\bigcup_iL_i)\cap\sphere^{N-1}$. Each constituent
is a sphere of dimension $N-a-1$; the intersection of one constituent
with any union of the others has dimension at most $N-b-1$. Applying
Mayer--Vietoris successively, and then Alexander duality, gives
\begin{equation}\label{eq:homologyrange}
 H_{a-1}(Z)\cong\Z^m,
 \qquad H_j(Z)=0\quad(a\le j\le b-3).
\end{equation}
The first isomorphism is given by small normal meridians about points
of $L_i$ lying on no other $L_j$. On the $i$th meridian the $i$th
component of $e$ has degree $1$, after choosing its orientation, and
every other component extends over the normal ball. Thus $e$ induces
an isomorphism in degree $a-1$.

The first possible product class in $Y$ above degree $a-1$ has degree
$2a-2$. The bound $b\le2a-1$ therefore gives $H_j(Y)=0$ for
$a\le j\le b-2$. The homology sequence of the mapping cylinder and
\eqref{eq:homologyrange} now show that its relative homology vanishes
through degree $b-2$. Both spaces are simply connected and have CW homotopy type. Successive
applications of the relative Hurewicz theorem give the corresponding
vanishing of relative homotopy. These are the usual duality and
Hurewicz arguments; see \cite{Hatcher}.
\end{proof}

\begin{remark}\label{rem:range}
The last degree matters. New homology of the source can occur in
degree $b-2$; the lemma does not assert its absence. It asserts a
relative connectivity statement, which is exactly what is needed to
lift maps of disks of dimension at most $b-2$ up to homotopy.
\end{remark}

\section{A scalar table in dimension six}\label{sec:scalar}

We prove Theorem~\ref{thm:A} by constructing a compatible map
$H:\Delta^5\to Z_{6,1}$. Write $I=[6]$, $V=V_{6,1}$ and $Z=Z_{6,1}$.
All matchings in this section are ordered by the label in their first
coordinate.

\subsection{Initial data}
Choose rational numbers $q_{i,a}$ and define the affine table
$H^0(\lambda)_w=\sum_i\lambda_iq_{i,w_i}$. For a matching $M$, put
$v_i=\ell_M(H^0(e_i))\in\R^5$. For a nonempty proper subset $A\subset I$
define the vector
\[
 u_A=6(\one_{1\in A}-\one_{6\in A},\ldots,
                   \one_{5\in A}-\one_{6\in A})\in\R^5.
\]
We require, for every $M$, that any five $v_i$ be linearly independent,
that $u_A$ and any four $v_i$ be independent when $1\le|A|\le3$, and
that $v_1,\ldots,v_6$ be affinely independent.

\begin{lemma}\label{lem:generic}
The numbers $q_{i,a}$ can be chosen to satisfy all these requirements.
One explicit choice is
\[
 B=65536,\qquad t_i=B^{7^{i-1}},\qquad q_{i,a}=t_i^a.
\]
\end{lemma}
\begin{proof}
For a permutation $\pi$ of $I$, the coefficient vectors of the
polynomial vector
$(t^{\pi(1)}-t^{\pi(6)},\ldots,t^{\pi(5)}-t^{\pi(6)})$
span $\R^5$: they are the five standard basis vectors and the negative
of their sum, in some order. Multilinearity makes every required
determinant, with a separate variable $t_i$ in each column, a nonzero
integer polynomial. The same argument applies to the affine determinant,
using the vectors $(1,v_i)$.

Each variable has degree at most six. The sum of absolute values of
coefficients is at most $23040$. For example the affine determinant
has bound $6\cdot5!\cdot2^5=23040$, and an anchored determinant has
bound $5!\cdot6\cdot2^4=11520$. Under the stated substitution distinct
monomials become distinct powers of $B$. The leading nonzero term
cannot be canceled by all lower terms, since $B>23040$.
\end{proof}

On a proper face of $\Delta^5$, at most five masses are positive.
The first independence condition therefore gives
$H^0(\partial\Delta^5)\subset Z$. This map already has the required
compatibility. We next modify its facets without changing their boundaries.

\subsection{Forty-one shells on a facet}
Let $J\subset I$ have five elements and let $\Delta_J$ be the corresponding
facet. Denote its barycenter by $b_J$, put $j_0=\min J$, and set
$q_J=H^0(b_J)$, $p_J=H^0(e_{j_0})$. For $A\subset I$ with
$1\le|A|\le3$ define the anchor table
$c_{J,A}(w)=6\one_{w_{j_0}\in A}-|A|$.

Define a map $T_{J,A}:[0,1]\times\partial\Delta_J\to V$ by
\begin{equation}\label{eq:anchor}
T_{J,A}(s,\theta)=
\begin{cases}
 (1-4s)q_J+4sp_J,&0\le s\le\tfrac14,\\
 (2-4s)p_J+(4s-1)c_{J,A},&\tfrac14\le s\le\tfrac12,\\
 (2-2s)c_{J,A}+(2s-1)H^0(\theta),&\tfrac12\le s\le1.
\end{cases}
\end{equation}
It starts at $q_J$, independently of $\theta$, and ends at $H^0(\theta)$.
Every value is safe. Indeed, after applying $\ell_M$, the first segment
is a nontrivial combination of five independent $v_i$ with nonnegative
coefficients. The second joins one $v_i$ to an anchor vector $u_P$,
where $P$ is a permutation of $A$. On the third segment, $H^0(\theta)$
uses at most four columns, which are independent together with $u_P$.

List the $21$ subsets of sizes one and two as $P_1,\ldots,P_{21}$,
and the $20$ subsets of size three as $N_1,\ldots,N_{20}$. Use radial
coordinates $\lambda=(1-r)b_J+r\theta$ on $\Delta_J$. Divide the radial
interval into $41$ equal subintervals and, from the center outwards,
use
\begin{equation}\label{eq:shells}
 T_{J,P_1},\ T_{J,N_1}^{\mathrm{rev}},\ T_{J,P_2},\
 T_{J,N_2}^{\mathrm{rev}},\ \ldots,\
 T_{J,P_{20}},\ T_{J,N_{20}}^{\mathrm{rev}},\ T_{J,P_{21}}.
\end{equation}
Here reversal means replacing $s$ by $1-s$. Adjacent shells agree
exactly: their common value is alternately $H^0(\theta)$ and $q_J$.
The center has value $q_J$ and the outer boundary has value $H^0(\theta)$.
The resulting facet maps glue to a continuous compatible safe map
$b:\partial\Delta^5\to Z$.

\subsection{The degree calculation}
Fix a matching $M$. For ordered vectors $z_0,\ldots,z_5\in\R^5$ such
that zero misses the affine images of proper faces, let
$D(z_0,\ldots,z_5)\in\Z$ be the degree at zero of the affine simplex,
with its ordered-vertex orientation. Put
\[
 o_M=D(v_1,\ldots,v_6),\qquad
 \phi_A(J)=D(u_A,v_{j_1},\ldots,v_{j_5})
 \quad(j_1<\cdots<j_5).
\]
The boundary of the abstract oriented simplex $[U,V_1,\ldots,V_6]$
gives, after affine mapping to $\R^5$, the local-homology identity
\begin{equation}\label{eq:degreeidentity}
 \sum_{i=1}^6(-1)^{i-1}\phi_A(I\setminus\{i\})=o_M.
\end{equation}
All terms are legitimate relative chains since their boundaries miss
zero. Their sum is a boundary, so its local degree is zero.

Replacing the original affine filling on an oriented facet by one
outward anchor filling changes its normal degree by $-\phi_P(J)$.
To see the sign, collapse the initial path in \eqref{eq:anchor}; this
is a homotopy in the complement of zero. The difference of the new
cone and the original facet is the negative of
$\partial[u_P,v_{j_1},\ldots,v_{j_5}]$. A reversed filling contributes
the opposite sign. This is an ordered-vertex computation; it does
not identify that orientation with a coordinate chart without its sign.

Write $\Phi_k(J)=\sum_{|A|=k}\phi_A(J)$. For a fixed matching the
selected color merely permutes the row labels, so summing over all
subsets of a given size removes this permutation. The interleaved
sequence \eqref{eq:shells} changes the facet class by
$-\Phi_1(J)-\Phi_2(J)+\Phi_3(J)$. One can compute this by collapsing
the spherical seams mapped constantly to $q_J$: each pair of an outward
and an inward shell then contributes the difference of its two fillings. There are $21$ outward and $20$
inward fillings, so the original facet is counted once. Including the
boundary sign of each facet and applying \eqref{eq:degreeidentity},
we obtain
\begin{equation}\label{eq:cancellation}
 \deg\left(\frac{\ell_M b}{\|\ell_M b\|}\right)
 =\left(1-\binom61-\binom62+\binom63\right)o_M=0.
\end{equation}
This holds for every matching, not just for a sampled collection.

By Lemmas~\ref{lem:codim} and~\ref{lem:comparison}, with $a=5,b=7$,
the normal-direction map from $Z$ to the product of $4$-spheres is
$5$-connected, and in particular an isomorphism on $\pi_4$.
Equation~\eqref{eq:cancellation} says that the image of $b$ represents
zero in the $\pi_4$ of every factor. Hence $b$ is null-homotopic in $Z$
and extends to $H:\Delta^5\to Z$. Its boundary is unchanged, so it
is compatible. Lemma~\ref{lem:realization} with $c=6,k=1$ proves
Theorem~\ref{thm:A}.

\section{Three height coordinates and twenty colors}\label{sec:vector}

The larger simplex required for Theorem~\ref{thm:concentrated} is obtained by a
relative lifting construction. We give the construction for a range
of parameters, since the dimension restriction is then transparent.

\begin{proposition}\label{prop:vector}
Let $k\ge2$ and $1\le c\le7k-1$ be integers. There is a continuous
compatible safe table $H:\Delta^{c-1}\to Z_{c,k}$, and consequently
an almost six-embedding $\Delta^{6c-1}\to\R^{c+k-1}$.
\end{proposition}

Let $W_6=\{(y_1,\ldots,y_6)\in\R^6:\sum_i y_i=0\}$, with the
permutation action. We use the following published input:
for a finite complex $X$ with a free $\Sym_6$-action there is an
equivariant map $X\to\sphere(W_6\oplus W_6)$
\cite[Theorem 2.2]{AKS}. The paper states the target as six planar
points not all equal; subtracting the mean and normalizing gives this
sphere. Since $k\ge2$, its inclusion in
$\sphere(W_6\otimes\R^k)$ is equivariant.

\subsection{A natural diagram of tests}
For each nonempty subset $S\subset[c]$, let
$\mathcal W_S=[6]^S$, $V_S=(\R^k)^{\mathcal W_S}$, and let $Z_S$ be
its safe-table space. The set of ordered matchings is
$\Omega_S=(\Sym_6)^S$: a tuple of permutations specifies the label
in each row of each active color. Relabeling the six rows gives a
free action on $\Omega_S$: explicitly $g\cdot(\pi_i)_i=(\pi_i g^{-1})_i$. Define
\[
 T_S=\operatorname{Map}_{\Sym_6}
       \bigl(\Omega_S,\sphere(W_6\otimes\R^k)\bigr).
\]
Since $\Omega_S$ is finite and free, $T_S$ is a finite product of
$(5k-1)$-spheres, one for each unordered matching.

Evaluating a table on an ordered matching, subtracting the mean of
its six vectors, and normalizing defines $e_S:Z_S\to T_S$.
Lemma~\ref{lem:comparison}, with $a=5k$ and $b=7k$, implies that
$e_S$ is $(7k-2)$-connected. For a single color there is only one
unordered matching and $e_S$ is a homotopy equivalence.

If $S\subset T$, duplicating a table independently of the extra labels
defines $j_{ST}:Z_S\to Z_T$. Precomposition with restriction of
ordered matchings defines $i_{ST}:T_S\to T_T$. These maps satisfy
\begin{equation}\label{eq:naturality}
                       e_Tj_{ST}=i_{ST}e_S
\end{equation}
exactly. A projected matching is still a matching, which also shows
that duplication preserves safety.

Consider the finite free complex $X=(\Sym_6)^{*c}$. Its points consist
of a mass vector and one permutation for each positive mass, with
inactive permutations forgotten. The published input supplies a map
$X\to\sphere(W_6\otimes\R^k)$. On every face $\Delta_S$ of the mass
simplex it defines a map $\beta_S:\Delta_S\to T_S$. On a smaller face
these maps agree after applying $i_{ST}$.

\subsection{Lifting, including the boundary homotopies}
We construct, by induction over faces, maps $H_S:\Delta_S\to Z_S$
and homotopies $e_SH_S\simeq\beta_S$, all compatible with inclusions
of faces. Suppose the construction has been made on the proper faces
of $\Delta_S$. Equation~\eqref{eq:naturality} gives a map of its
boundary to $Z_S$ and a specified homotopy there to $\beta_S$.

The relative lifting fact being used is this: if $e:A\to B$ is an
$n$-connected map of spaces of CW homotopy type and a disk has dimension
at most $n$, a map of the disk
to $B$ lifts up to homotopy from prescribed boundary lift and boundary
homotopy. Indeed replace $e$ by its mapping-path fibration. Its
homotopy fiber is $(n-1)$-connected; the obstruction to extending a
section over a $d$-cell belongs to $\pi_{d-1}$ of that fiber and
vanishes for $d\le n$. Projection to $A$ recovers the desired lift
and its homotopy, with the prescribed boundary data unchanged.

Here $\dim\Delta_S\le c-1\le7k-2$. The lifting fact therefore applies
to $e_S$ at every stage, including the largest face. The resulting map
$H=H_{[c]}$ is continuous, safe, and compatible. This proves
Proposition~\ref{prop:vector}. Taking $k=3,c=20$ gives
$6c-1=119$ and $c+k-1=22$, and proves Theorem~\ref{thm:C}.

\begin{remark}[A second proof of Theorem~\ref{thm:A}]
The same diagram gives an independent route to the scalar existence
statement. The \"Ozaydin theorem, as stated in \cite[Theorem 2.2$'$]{AKS},
gives an equivariant map from the five-dimensional free complex
$(\Sym_6)^{*6}$ to $\sphere(W_6)$. For $k=1$, each comparison $e_S$
is $5$-connected by Lemma~\ref{lem:comparison}. The facewise induction
therefore lifts this test over $\Delta^5$, and Lemma~\ref{lem:realization}
gives Theorem~\ref{thm:A}. This proof shares the arrangement-comparison
lemma but does not use the shell signs or the degree cancellation.
The explicit scalar construction above supplies the rational boundary
prescription absent from this shorter existence argument.
\end{remark}

\begin{remark}
The homotopies are part of the induction, not an optional addition.
An arbitrary equivariant test is not automatically a height table.
The connectivity of $e_S$ and the relative lifting argument supply
precisely this missing realization step, while retaining compatibility
on every zero-mass face.
\end{remark}

\subsection{Padding by five vertices}
The following elementary join construction is the special case of
\cite[Remark 1.2(a)]{AKS} in which the second factor has only five
vertices. Its proof is included to fix the dimension count.
\begin{lemma}\label{lem:padding}
If there is an almost six-embedding $g:\Delta^N\to\R^q$, there is an
almost six-embedding $g^+:\Delta^{N+5}\to\R^{q+1}$.
\end{lemma}
\begin{proof}
Regard $\Delta^{N+5}$ as the join $\Delta^N*\Delta^4$. Write a join
point as $t x*(1-t)y$, with $t\in[0,1]$, $x\in\Delta^N$, and
$y\in\Delta^4$, and set
\[
                    g^+(t x*(1-t)y)=(t,tg(x)).
\]
The formula is continuous, including at $t=0$, since $g$ is bounded.
In a sixfold coincidence the first coordinates force a common value
of $t$. If $t>0$, the six old components have disjoint supports and
a common $g$-image, which is impossible. If $t=0$, all six supports
would lie among the five new vertices, which is also impossible.
\end{proof}

\begin{corollary}\label{cor:wideall}
For every integer $D\ge24$ there is an almost six-embedding
\[
                       g_D:\Delta^{5D-1}\longrightarrow\R^{D-2}.
\]
\end{corollary}
\begin{proof}
Starting from Theorem~\ref{thm:C}, apply Lemma~\ref{lem:padding}
$D-24$ times. The source dimension becomes
$119+5(D-24)=5D-1$, and the target dimension becomes
$22+(D-24)=D-2$.
\end{proof}

\section{How many parameters does a six-cone partition need?}\label{sec:cones}

A \emph{central cone partition} of $\R^D$ is a finite tuple of closed,
full-dimensional convex polyhedral cones covering $\R^D$ with disjoint
interiors. Its internal skeleton is the complement of the union of
the interiors. No face-to-face condition is imposed.

A dimension bound for this family means that finitely many semialgebraic
parameter spaces of the stated dimensions represent all its skeleta.
Repeated representations cause no difficulty. Unit normals remove
positive rescaling; all incidence, coverage, and full-dimensionality
conditions are semialgebraic by real quantifier elimination; the dimension and projection facts used below are standard, as in \cite{BPR}.

\begin{theorem}\label{thm:conedim}
For $D\ge4$, the internal skeleta of central cone partitions of $\R^D$
into at most six cones have a finite semialgebraic parametrization of
dimension at most
\begin{equation}\label{eq:conebound}
                         \max\{5D-1,\,4D+23\}.
\end{equation}
\end{theorem}

\subsection{The effective normal space}
For a partition $K_1,\ldots,K_r$, choose nonzero separating normals
$n_{ij}=-n_{ji}$ so that $K_i\subset\{x:\langle n_{ij},x\rangle\ge0\}$.
Then
\begin{equation}\label{eq:pairseparators}
 K_i=\bigcap_{j\ne i}\{x:\langle n_{ij},x\rangle\ge0\}.
\end{equation}
For if the right side contained a point outside $K_i$, a segment to an
interior point of $K_i$ would give an open set of such points satisfying
all inequalities strictly. An interior point of another cone in this
open set would violate its pairwise separator with $K_i$.

\begin{lemma}[Normal rank]\label{lem:rank}
The span $F$ of all pair-separating normals has dimension at most $r-1$.
It equals the span of the genuine facet normals of all cells.
\end{lemma}
\begin{proof}
Put $a_i=\sum_{j\ne i}n_{ij}$, so $\sum_i a_i=0$. We prove that the
$a_i$ span $F$. Suppose a vector $v$ annihilates all $a_i$, and let
$S=\{i:v\in K_i\}$. This set is nonempty. For $i\in S$, all
$\langle n_{ij},v\rangle$ are nonnegative and sum to zero, so all vanish.
If $k\notin S$, take $x\in\intt K_k$. For $i\in S$,
$\langle n_{ik},x+tv\rangle<0$ for every $t$. For $i\notin S$, one
of its defining normals is negative on $v$, so $x+tv\notin K_i$ for
all large $t$. This contradicts coverage. Thus $S=[r]$ and $v$
annihilates every normal. The span assertion follows by orthogonal
complements. Genuine facets occur among the defining hyperplanes in
\eqref{eq:pairseparators}; conversely every valid defining normal lies
in the dual cone, hence in the span of the genuine facet normals.
\end{proof}

We call $q=\dim F$ the normal rank. The common lineality of the
partition is $F^\perp$. Choosing $F$ costs $q(D-q)$ parameters.
For $1\le q\le3$, choosing all fifteen pairwise directions inside $F$ gives
at most $q(D-q)+15(q-1)$ parameters, which is at most $3D+21$ for
$D\ge4$. We need sharper arguments in ranks four and five.

\subsection{The maximal-rank case}
\begin{lemma}\label{lem:maxrank}
Let $r\ge2$ be an integer. Partitions into $r$ cones with normal rank $r-1$ have a semialgebraic
parametrization of dimension at most $(r-1)D-1$.
\end{lemma}
\begin{proof}
Work first in the effective $(r-1)$-space. The map
$x\mapsto(\langle a_i,x\rangle)_i$ identifies it with
$W_r=\{y\in\R^r:\sum_i y_i=0\}$. Each $K_i$ lies in $y_i\ge0$.
Points with $y_i>0$ and every other coordinate negative belong to
$\intt K_i$.

For three indices $i,j,k$, fix all other coordinates at strictly
negative values. The resulting affine plane is covered by exactly
the three possible cells. Each has a two-dimensional private recession
sector. The map from this plane defined by the three cyclic separators
$(n_{ij},n_{jk},n_{ki})$ has rank two: rank at most one would give
three disjoint intervals along one coordinate, each with a
full-dimensional recession sector, whereas at most two disjoint
intervals can be unbounded.

Its affine-plane image in $\R^3$ misses both strictly positive and
strictly negative orthants; either sign pattern leaves no cell able
to contain the point. An affine plane missing both orthants passes
through zero. Indeed its equation $c\cdot z=b$ meets both orthants
if $c$ has both signs; otherwise it meets one of them unless $b=0$.
As the fixed negative coordinates vary, the common kernel of the
three global separators projects onto an open subset of $\R^{r-3}$.
Their rank is therefore at most two. Thus every triangle of pairwise
normal directions is linearly dependent.

We record the dimension count for these directions. Set $q=r-1$.
Assign a nonzero projective direction to each edge of the complete
graph on $q+1$ vertices, assume that they span $\R^q$, and impose
dependence on every triangle. Their parameter space has dimension at
most $q^2-1$. To prove this, first note that independent edges form a
forest: triangulating a hypothetical independent cycle successively
puts each nonzero chord in the span of its preceding path, contradicting
the last triangle. In particular directions on $m$ vertices span at
most $m-1$ dimensions.

Proceed by induction on $q$. If all vertex stars have rank $q$, choose
the $q$ independent projective directions at one vertex, represented by
vectors $w_1,\ldots,w_q\in\R^q$. This costs $q(q-1)$ parameters. Every other edge has direction
$[t_iw_i-t_jw_j]$ with all $t_i\ne0$; this follows from the triangle
relations and their multiplicatively consistent coefficient ratios.
The ratios have $q-1$ parameters, giving $q^2-1$.

Otherwise choose a deficient star. Group its neighbors into $s$ classes
$C$ of parallel directions, let $U$ be the star span of dimension $k$,
and let $V_C$ be the span of the edges within $C$. Cross-class edges
belong to $U$, and $\dim V_C\le |C|-1$. Consequently
$q\le k+\sum_C\dim V_C\le k+q-s\le q$.
Equality forces $k=s$, $\dim V_C=q_C=|C|-1$, and the direct sum
$\R^q=U\oplus\bigoplus_{q_C>0}V_C$.
A triangle with two vertices in one class shows that its cross-class
edges are parallel, so they descend to a configuration on the classes
and the original vertex. Apply induction to that configuration and
to each nontrivial class. The choice of the ordered direct sum has
dimension $q^2-s^2-\sum_Cq_C^2$; adding the internal bounds gives
at most $q^2-1$. There are only finitely many class patterns.

Finally choose the effective space in $q(D-q)$ parameters. The total
is $q(D-q)+q^2-1=qD-1$. Orienting projective directions is a finite
choice, and the cones are recovered from their separators.
\end{proof}

\subsection{Ridge constraints in rank four}
We use a local fact which does not require face-to-face intersections.

\begin{lemma}[A ridge needs another wall]\label{lem:ridge}
Let $K_0$ be a cell and let $P$ be the two-dimensional normal space
of one of its genuine codimension-two faces. At least one genuine
interface between cells other than $K_0$ has its normal in $P$.
\end{lemma}
\begin{proof}
Choose a point in the relative interior of the face, avoiding all
other wall hyperplanes which do not contain its linear span. In the
normal two-plane the local section of $K_0$ is a pointed convex angle,
strictly smaller than $\pi$. Its complementary angle is larger than
$\pi$ and cannot belong to one convex cell. There is therefore an
interface between two other local cells. Every wall through the chosen
point contains the face span, by the generic choice, so the interface
normal lies in $P$. The local interface has a relatively open
codimension-one patch, hence is a genuine global interface.
\end{proof}

\begin{lemma}\label{lem:rankfour}
Six-cone partitions of normal rank four in $\R^D$, $D\ge4$, have a
semialgebraic parametrization of dimension at most $4D+23$.
\end{lemma}
\begin{proof}
Fix a labeled reference cell $K_0$, let $h$ be its normal rank, and
let $t\le5$ be its number of genuine facets. Choose its normal space
$E$, its $t$ unit facet normals, and then the total four-space $F$
containing $E$. These choices cost
$h(D-h)+t(h-1)+(4-h)(D-4)$ parameters.
The ten pairwise separator directions among the other five cells
cost at most $30$ more. The initial bound is thus
\begin{equation}\label{eq:rankfourbase}
                  4D+14+h(4-h)+t(h-1).
\end{equation}
These data determine the partition: for $i\ne0$, the cell $K_i$ is
the closure of the points outside $K_0$ that satisfy strictly all its
pairwise inequalities against the other four cells. Coverage proves
one inclusion, and $\intt K_i$ proves the other. Thus no extra
separators incident with $K_0$ need to be counted.

If $h=1$, then $t=1$ and \eqref{eq:rankfourbase} is $4D+17$.
If $h=2$, then $t=2$ and it is $4D+20$.

Suppose $h=3$. A section of the pointed dual cone of $K_0$ is a polygon
with $t$ vertices. If $t=3$, the bound is already $4D+23$. If $t=4$,
one ridge, by Lemma~\ref{lem:ridge}, puts one internal normal in a
specified two-plane of $F$. This costs two parameters in projective
three-space and reduces $4D+25$ to $4D+23$.
If $t=5$, take three consecutive edge planes of the dual pentagon.
Their common intersection is zero: the first two intersect in the
second vertex ray, which is not in the plane of the third edge.
Each plane must contain an internal separator. If three distinct
separator slots witness them, six parameters are lost. If two planes
use the same slot, that direction is confined to their line of
intersection, costing three parameters, and the third plane costs
two more. One slot cannot serve all three. The loss is at least five,
reducing $4D+27$ to $4D+22$.

Suppose $h=4$. A section of the dual cone is a three-dimensional
polytope with four or five vertices. Its graph contains a four-cycle
on four affinely independent vertices. For a tetrahedron this is
immediate. A three-polytope with five vertices has eight or nine edges, by
Euler's formula and minimum vertex degree three. With eight edges its
graph is a four-cycle joined to an apex; with nine it is $K_5$ with
one edge removed. Thus it is a square pyramid or a triangular
bipyramid. In the pyramid use the apex and three consecutive base
vertices. In the bipyramid use both apices and two equatorial vertices;
no four vertices are coplanar, since they would form a quadrilateral
facet. In either case the selected four-cycle is affinely independent.
Choose nonzero vectors $u_1,u_2,u_3,u_4$ on the corresponding dual
rays, in cyclic order. These four vectors are linearly independent. The four ridge normal planes are
\[
 P_{12}=\spanlin(u_1,u_2),\quad P_{23}=\spanlin(u_2,u_3),\quad
 P_{34}=\spanlin(u_3,u_4),\quad P_{41}=\spanlin(u_4,u_1).
\]
Opposite planes intersect only at zero; adjacent ones intersect in
a line. Lemma~\ref{lem:ridge} assigns an internal separator to each
plane. A separator can serve at most two planes, and only adjacent
ones. Four distinct slots cost eight parameters; one shared slot
and two single slots cost seven; two shared slots cost six. The
loss is therefore at least six. Since $t\le5$, the bound
$4D+14+3t$ falls to at most $4D+23$.

All assignments of facet labels and witnessing separator slots form
finite lists. The losses were computed in fibers over the reference
cell and total normal space, so no independence of varying reference
planes was assumed. Taking their finite union proves the bound.
\end{proof}

\begin{proof}[Proof of Theorem~\ref{thm:conedim}]
For six cells, Lemma~\ref{lem:rank} gives rank at most five.
Lemma~\ref{lem:maxrank} gives $5D-1$ in rank five;
Lemma~\ref{lem:rankfour} gives $4D+23$ in rank four; the bound
$3D+21$ in lower ranks is smaller when $D\ge4$.
For five cells, rank four is maximal and costs at most $4D-1$ by
Lemma~\ref{lem:maxrank}; all remaining cases with two to five cells
have rank at most three and obey the crude bound already given.
The one-cell skeleton is empty. This proves the parametrization
bound for every number of cells at most six. A smaller partition can
also be refined to six by central cuts without losing any old
internal boundary.
\end{proof}

\section{From blocking points to body functionals}\label{sec:transfer}

\begin{lemma}[Generic blocking points]\label{lem:blockers}
For every integer $D\ge24$, there are $5D$ distinct rational points
$p_1,\ldots,p_{5D}\in\intt B_D$ which do not all lie on the internal
skeleton of any central cone partition of $\R^D$ into at most six
full-dimensional cells.
\end{lemma}
\begin{proof}
For $D\ge24$, Theorem~\ref{thm:conedim} gives parameter dimension
$\max\{5D-1,4D+23\}=5D-1$. For a fixed skeleton each point has at
most $D-1$ dimensions of choice. The incidence set of a skeleton and
$5D$ points on it therefore has dimension at most
\[
                (5D-1)+5D(D-1)=5D^2-1<5D^2.
\]
Its projection, and the closure of that projection, have smaller
dimension than the space of point tuples. Their complement contains
a nonempty open subset of $(\intt B_D)^{5D}$. Choose a rational tuple
with distinct entries there. All smaller partitions were included in
Theorem~\ref{thm:conedim}. For $D=24$ the two incidence dimensions are
$2879$ and $2880$, and there are $120$ points.
\end{proof}

\subsection{The compact geometric obstruction}
For this subsection only, a \emph{generalized six-partition} of $B_D$
is a tuple of six nonempty compact convex sets covering $B_D$ with
pairwise disjoint interiors; empty interiors are allowed. These
tuples form a compact space $\mathcal Q_D$. Indeed the hyperspace of
nonempty compact convex subsets of $B_D$ is compact. Coverage persists
under Hausdorff limits. Disjointness of interiors also persists,
because a compact ball in the interior of a full-dimensional convex
limit is contained in its approximants for all sufficiently large
indices.

\begin{lemma}[A positive distance from centrality]\label{lem:nearcentral}
Let $D,m\ge1$ be integers, and let $p_1,\ldots,p_m\in\intt B_D$
block all central cone partitions with at most six cells. There is
$\eta\in(0,1)$ such that every generalized six-partition in which
none of the marked points is interior to any part has a part $C_i$
with $\dist(0,C_i)>\eta$.
\end{lemma}
\begin{proof}
In $\mathcal Q_D$, the condition that all inward distances
$\dist(p_a,\R^D\setminus C_i)$ vanish is closed. On the resulting
compact set the function $\max_i\dist(0,C_i)$ is continuous. It
suffices to prove that it never vanishes; if the set is empty, the
claim is vacuous.

Suppose it vanishes at a tuple. Every part then contains zero.
Discard the parts with empty interior. The remaining sets still
cover $B_D$: their finite union is closed, and an uncovered relatively
open portion of the ball could not be covered by finitely many
lower-dimensional convex sets. Choose separating hyperplanes for
pairs of surviving parts. Because both parts contain zero, each
separator passes through zero. Intersect the corresponding halfspaces
to obtain cones $K_i$. Coverage, as in \eqref{eq:pairseparators}, gives
$C_i=B_D\cap K_i$. Every direction has a short representative in the
ball, so the cones cover $\R^D$ and have disjoint interiors. All marked
points lie on their internal skeleton, contradicting the blocking
property. The continuous function consequently has a positive minimum;
choose $\eta<1$ smaller than that minimum.
\end{proof}

Allowing degenerate parts here is deliberate. It gives a compact
closed set on which a positive separation margin can be taken; no
measure-convergence assertion is needed for this step.

\begin{lemma}[Inward-distance transfer with a uniform margin]\label{lem:transfer}
Let $D,m,q$ be positive integers. Suppose that
$g:\Delta^{m-1}\to\R^q$ is an almost six-embedding and that
$m$ points in $\intt B_D$ block all central cone partitions with
at most six cells. There exist a continuous map
$\Phi:\mathcal K_D\to\R^{q+1}$ and numbers $\eta\in(0,1)$ and
$\varepsilon>0$ such that every generalized six-partition whose
parts all meet $\eta B_D$ satisfies
\begin{equation}\label{eq:margin}
                  \max_{1\le i\le6}\|\Phi(C_i)-\Phi(C_1)\|_\infty
                     \ge\varepsilon.
\end{equation}
\end{lemma}
\begin{proof}
Let the blocking points be $p_1,\ldots,p_m$. For $C\in\mathcal K_D$
define its inward distances $w_a(C)=\dist(p_a,\R^D\setminus C)$,
and put $s(C)=\sum_a w_a(C)$. Set
\begin{equation}\label{eq:functionals}
 \Phi(C)=
 \begin{cases}
 \bigl(s(C),\,s(C)g(w(C)/s(C))\bigr),&s(C)>0,\\
 0,&s(C)=0.
 \end{cases}
\end{equation}
This definition includes lower-dimensional sets. If $h_C$ is the
support function of $C$, then
\begin{equation}\label{eq:inward}
 w_a(C)=\max\left\{0,\min_{\|u\|=1}
                    \bigl(h_C(u)-\langle u,p_a\rangle\bigr)\right\}.
\end{equation}
Hence $w_a$ is $1$-Lipschitz in the Hausdorff metric. Boundedness of
$g$ on its compact domain proves continuity of \eqref{eq:functionals}
at $s=0$; elsewhere continuity is immediate.

Suppose a generalized six-partition equalizes $\Phi$. If the common
value of $s$ is $t>0$, the six points $w(C_i)/t$ of
$\Delta^{m-1}$ have disjoint supports: a marker cannot be interior
to two parts. Equality of the other coordinates contradicts the
almost six-embedding property. Thus equality forces
\begin{equation}\label{eq:zero}
                         w_a(C_i)=0\quad\text{for all }a,i.
\end{equation}
Choose $\eta$ from Lemma~\ref{lem:nearcentral}. Equation~\eqref{eq:zero}
is impossible when all parts meet $\eta B_D$. The set of such tuples
in $\mathcal Q_D$ is compact, since distance from zero is continuous.
The contrast in \eqref{eq:margin} is continuous and nowhere zero on
this set, so it has a positive minimum. If the set is empty, any
positive $\varepsilon$ suffices.
\end{proof}

\subsection{One measure, with the quantifiers in the correct order}
The following elementary observation is the only measure-theoretic
step of the transfer.
\begin{lemma}\label{lem:mass}
Let $0<\eta<1$ and let $\mu$ be a Borel measure with
$0<\mu(B_D)<\infty$. If
$\mu(\eta B_D)>5\mu(B_D)/6$, then every $\mu$-fair six-partition
of $B_D$ has each of its parts meeting $\eta B_D$.
\end{lemma}
\begin{proof}
A part $C_i$ missing $\eta B_D$ is contained in
$B_D\setminus\eta B_D$, and hence
\[
 \mu(C_i)\le\mu(B_D)-\mu(\eta B_D)<\frac{\mu(B_D)}6,
\]
contrary to fairness. Positivity and finiteness of $\mu(B_D)$ justify
this comparison. Absolute continuity is not needed in this lemma.
\end{proof}

\begin{proof}[Proof of Theorem~\ref{thm:concentrated} and Theorem~\ref{thm:B}]
Fix $D\ge24$. Use $g_D:\Delta^{5D-1}\to\R^{D-2}$ from
Corollary~\ref{cor:wideall} and the $5D$ points of
Lemma~\ref{lem:blockers}. Lemma~\ref{lem:transfer} gives a fixed map
$\Phi_D$ with $(D-2)+1=D-1$ coordinates and a fixed radius $\eta_D$.
For every measure in the statement, Lemma~\ref{lem:mass} puts every
fair partition in the compact class covered by \eqref{eq:margin}.
It cannot equalize $\Phi_D$. Neither the map nor the radius was
chosen using the measure. This proves Theorem~\ref{thm:concentrated}.

For Theorem~\ref{thm:B}, take $d=D\ge24$, the map $\Phi_d$, and any
$0<\sigma<\eta_d/\sqrt{6d}$. As shown in Corollary~\ref{cor:gaussian}
below, the centered Gaussian $\gamma_{d,\sigma}$ satisfies
$\gamma_{d,\sigma}(\eta_d B_d)>(5/6)\gamma_{d,\sigma}(B_d)$, so its
normalized restriction $\mu(E)=\gamma_{d,\sigma}(E\cap B_d)/\gamma_{d,\sigma}(B_d)$
is an absolutely continuous probability measure on $B_d$ meeting the
hypotheses of Theorem~\ref{thm:concentrated}.
\end{proof}

The first coordinate $s(C)$ in \eqref{eq:functionals} is a body
functional, not a measure: it is not additive on partitions. It
is one of the $D-1$ coordinates already counted above. The single
measure $\mu$ supplies a separate equal-mass condition, which forces
all parts to meet a common small ball. This is why no additional
centrality coordinate is required.

\begin{proof}[Proof of Corollary~\ref{cor:gaussian}]
Let $X$ be a random vector with distribution $\gamma_{D,\sigma}$.
Since $\mathbb E\|X\|^2=D\sigma^2$, Markov's inequality gives
\[
 \gamma_{D,\sigma}(\eta_D B_D)
 \ge 1-\frac{D\sigma^2}{\eta_D^2}.
\]
For $0<\sigma<\eta_D/\sqrt{6D}$ this is strictly greater than
$5/6$, and therefore greater than $5\gamma_{D,\sigma}(B_D)/6$.
Apply Theorem~\ref{thm:concentrated}. For $D=24$, $\sqrt{6D}=12$.
The Gaussian density is everywhere positive and analytic, and the
same functionals work throughout the stated interval of scales.
\end{proof}

More generally, the conclusion holds for any absolutely continuous
probability measure $\mu$ on $\R^D$ with $\int\|x\|^2\dd\mu(x)<\eta_D^2/6$.
The concentration criterion is invariant under multiplication of a
measure by a positive constant. It needs neither radial symmetry nor
a Gaussian density. Requiring equality of a further measure would
only reduce the set of admissible partitions, but it would be a
different problem, with two measures instead of one.

\begin{corollary}[A prescribed measure on a variable ball]\label{cor:prescribed}
Let $D\ge24$ and let $\mu$ be any finite nonzero absolutely continuous
Borel measure on $\R^D$. There exist $R>0$ and $D-1$ continuous
body functionals which cannot be equalized, together with $\mu$,
on a convex six-partition of $RB_D$.
\end{corollary}
\begin{proof}
Let $M=\mu(\R^D)\in(0,\infty)$ and let $\eta_D$ be as above.
Choose $R>0$ so large that $\mu(\eta_D RB_D)>5M/6$. Then
$\mu(\eta_D RB_D)>5\mu(RB_D)/6$. Under the dilation $x\mapsto x/R$,
the pushforward measure meets the hypotheses of
Theorem~\ref{thm:concentrated}. Apply that theorem with
$\Phi_R(C)=\Phi_D(R^{-1}C)$. Dilation is continuous in the
Hausdorff metric and preserves convexity, coverage, and disjoint
interiors. This proves the assertion.
\end{proof}
This corollary allows the ball to depend on the prescribed measure.
It does not convert the theorem into a statement for every prescribed
measure on a fixed ball. In particular, it gives the standard Gaussian
on a sufficiently large ball, not a new claim for the standard Gaussian
on the unit ball.

\subsection{Moment polynomials and robustness}
\begin{proposition}[A finite-moment version]\label{prop:moments}
The map $\Phi_D$ in Theorem~\ref{thm:concentrated} may be chosen with coordinates
that are polynomials with rational coefficients in the ordinary
Lebesgue moments $m_\alpha(C)=\int_Cx^\alpha\dd x$, $|\alpha|\le12$.
One such map works for every measure in the stated concentration
class, with the same $\eta_D$. No bound on its degree as a polynomial
in the moments is asserted.
\end{proposition}
\begin{proof}
Every part of a convex six-partition of $B_D$ is the intersection
of the ball with at most five halfspaces, by pairwise separation
and coverage. The same is true of any full-dimensional Hausdorff
limit of such parts: pass to convergent unit separator normals and
offsets, and test strict inequalities in the interior of the
limiting intersection. Let $\mathcal A$ be the closure of the family
of all such parts in $\mathcal K_D$. It is compact, and its other
members are lower-dimensional.

Moments through order $12$ distinguish its full-dimensional members.
Suppose that $C,K\in\mathcal A$ are full-dimensional and have equal
moments. Let $\ell_1=0,\ldots,\ell_t=0$, $t\le5$, be the genuine
flat facet planes of $C$. For $z\in\intt K$ use the polynomial
vector field
\[
 V_z(x)=(x-z)(1-\|x\|^2)\prod_{j=1}^t\ell_j(x)^2.
\]
Its divergence has degree at most $2t+2\le12$, and its flux through
$\partial C$ is zero. Equality of moments therefore gives zero
flux through $\partial K$. On each flat facet of $K$ the outward
normal has positive scalar product with $x-z$, and all remaining
factors are nonnegative. The relatively open part of the facet
inside the ball must consequently lie in one of the planes of $C$.
A finite union of different hyperplanes cannot cover an open part
of one hyperplane. The common centroid lies in both interiors and
chooses the same halfspace sides. It follows that $C\subset K$,
and equal volumes give $C=K$.

Let $M$ be the vector of moments through order $12$. It is continuous
on $\mathcal A$. Its only possible nontrivial fibers consist of
lower-dimensional sets, on which all moments and the original
functional \eqref{eq:functionals} vanish. Thus that functional
descends continuously to $M(\mathcal A)$: a continuous map from a
compact space onto a Hausdorff space is a quotient map.

Use Stone--Weierstrass to approximate its coordinates uniformly
on $M(\mathcal A)$ by polynomials, and then approximate the finitely
many coefficients by rationals. Choose the coordinatewise error
smaller than $\varepsilon/4$, where $\varepsilon$ is the geometric
margin of Lemma~\ref{lem:transfer}. The contrast on every relevant
partition changes by less than $\varepsilon/2$ and stays positive.
Lemma~\ref{lem:mass} applies to every measure in the same class, so
no measure-dependent compactness or approximation is required.
The resulting polynomials in moments are continuous on all of
$\mathcal K_D$.
\end{proof}

The same margin shows that sufficiently small uniform perturbations
of these functionals on the compact family $\mathcal A$ preserve
the obstruction for the whole concentration class. This is robustness
of a counterexample, not an assertion that every collection of
$D-1$ functions is a counterexample.

\begin{remark}[The numerical threshold]
For a base construction with $c$ colors and $k$ height coordinates,
the transfer uses $6c$ markers and $c+k$ functionals, so its ambient
dimension is $D=c+k+1$. The sufficient inequalities are
\[
 c\le7k-1,\qquad c\ge5k+5,\qquad c\ge2k+14.
\]
The first is the realization range; the last two express
$6c>\max\{5D-1,4D+23\}$. For $k=2$ the first two conflict. For
$k\ge3$, $D\ge6k+6\ge24$, with equality at $k=3,c=20$.
Padding by five vertices propagates this base construction to every
$D\ge24$. Neither this optimization nor the padding is a
nonexistence argument below dimension $24$. The revision changes
the class of measures and the uniformity, not the lowest dimension
established by this argument.
\end{remark}

\section{Constructivity and the formal statements}\label{sec:formal}

The scalar construction specifies its boundary table by rational data
and a finite shell prescription. Before a rational piecewise-linear
search, replace this boundary by a rational PL map: triangulate the
radial shell prisms, with breaks at the three stages of
\eqref{eq:anchor}, and interpolate their vertex values. On each
small simplex both the original and interpolated images lie in one
of the safe convex sets used in the safety proof: a bridge segment,
or $\conv\{c_{J,A},H^0(e_i):i\in K\}$ with $|K|\le4$.
Straight interpolation gives a safe homotopy fixed on common facet
boundaries. Thus compatibility and all degrees are unchanged.

A safe extension of this PL boundary can be found in principle by
a fair enumeration of rational interior vertex data on successive
barycentric subdivisions, with exact linear-feasibility tests against
every $L_M$. Compactness supplies a positive avoidance margin for
an existing continuous extension; sufficiently fine rational
interpolation preserves it and occurs in the enumeration. This
proves termination, but gives no practical complexity bound or
expanded coefficient list.

The same type of finite search applies to the vector table, without
prescribing its boundary in advance. On successive rational
subdivisions of the mass simplex enumerate rational vertex tables
which ignore every inactive color. Affine interpolation then preserves
compatibility. Every candidate has a finite exact safety test against
all $L_M$. Uniform approximation of the compatible continuous table
from Proposition~\ref{prop:vector} proves that this search terminates.

The blocking points may also be specified by a terminating rational
search. For rational input points, existence of a central cone
partition with all points on its skeleton is a first-order statement
over the real field: use unit pairwise separator normals and
polynomial inequalities for fullness, coverage, and interior
membership. Real quantifier elimination decides it. Lemma~\ref{lem:blockers}
supplies an open set of blocking tuples, so enumeration of rational
tuples eventually finds one.

For a rational blocking tuple, one can likewise enumerate dyadic
$\eta>0$ and test whether there is a partition into $r$ full-dimensional
parts, for any $1\le r\le6$, having no interior marker and with every
part meeting $\eta B_D$. Each such part is the intersection of $B_D$
with at most $r-1$ affine halfspaces, so this is a finite first-order
test over the real field. This test also covers the generalized tuples
in Lemma~\ref{lem:nearcentral}: discard their empty-interior parts;
the remaining parts still cover the ball and retain all the tested
conditions. Conversely, a partition with fewer than six parts can be
padded with copies of $\{0\}$. Thus quantifier elimination verifies
the required exclusion without parametrizing arbitrary degenerate
convex sets. The lemma guarantees termination. No Gaussian integral
enters this test.
A rational Gaussian scale can then be chosen below
$\eta/\sqrt{6D}$. In dimension $D=24$, the interval endpoint
$\sigma_0=\eta/12$ is rational when $\eta$ is rational. Every scale
strictly below it has exterior mass less than $1/6$, so the inner
ball carries more than $5/6$ of the full Gaussian probability.
For Theorem~\ref{thm:concentrated}, the geometric data are fixed first and the
measure is tested only afterwards by its concentration inequality.

These are algorithms with existence-based termination proofs, not
executed coefficient computations or practical complexity bounds.
Once their outputs are fixed, \eqref{eq:functionals} is an exact formula
for the body functionals. Proposition~\ref{prop:moments} gives a
finite-moment alternative, but its coefficients have not been computed.
The proof never asserts that limiting cones have equal solid angles;
only common containment of zero and absence of interior markers are used.

\begin{sloppypar}
The accompanying Lean~4 formalization (pinned to \texttt{leanprover/lean4:v4.35.0}
and Mathlib~\texttt{v4.35.0}, configured for Palomar via
\texttt{Challenge.lean}, \texttt{Solution.lean}, \texttt{comparator.json},
and \texttt{formalization.yaml}) builds cleanly with zero \texttt{sorry}
in \texttt{Solution.lean} and all its dependencies.
Theorem~\ref{thm:A} (\texttt{AkopyanSixfold.\allowbreak tverberg\_\allowbreak counterexample})
is proved by executing the scalar height-table reduction, rational shell
degree balance, and piecewise-linear zero-cancellation surgery in
\texttt{RequestProject/}. Theorem~\ref{thm:B}
(\texttt{AkopyanSixfold.\allowbreak convex\_\allowbreak partition\_\allowbreak counterexample})
is proved in every dimension $d\ge24$ using standard Euclidean topologies, Mathlib's
Hausdorff metric on \texttt{ConvexBody}, and a Borel probability measure
(\texttt{IsProbabilityMeasure}), without restricting to absolutely
continuous measures. Kernel axiom inspection confirms that
both main theorems depend exclusively on \texttt{propext},
\texttt{Quot.sound}, and \texttt{Classical.choice}.
\end{sloppypar}

\subsection*{Questions left open by these proofs}
These arguments do not decide whether an almost six-embedding
$\Delta^{30}\to\R^5$ exists, or whether fair-partition counterexamples
exist in dimension $8$ or $7$. They give existence bounds, not
nonexistence theorems in smaller dimensions.

\begin{thebibliography}{12}
\bibitem{AAK}
A.~Akopyan, S.~Avvakumov, and R.~Karasev,
\emph{Convex fair partitions into an arbitrary number of pieces},
Adv. Math. \textbf{493} (2026), 110927.
\href{https://doi.org/10.1016/j.aim.2026.110927}{doi:10.1016/j.aim.2026.110927};
arXiv:1804.03057v14.

\bibitem{AKS}
S.~Avvakumov, R.~Karasev, and A.~Skopenkov,
\emph{Stronger counterexamples to the topological Tverberg conjecture},
Combinatorica \textbf{43} (2023), 717--727.
\href{https://doi.org/10.1007/s00493-023-00031-w}{doi:10.1007/s00493-023-00031-w};
arXiv:1908.08731v4.

\bibitem{AMS}
S.~Avvakumov, I.~Mabillard, A.~Skopenkov, and U.~Wagner,
\emph{Eliminating higher-multiplicity intersections, III. Codimension 2},
Israel J. Math. \textbf{245} (2021), 501--534.
\href{https://doi.org/10.1007/s11856-021-2216-z}{doi:10.1007/s11856-021-2216-z}.

\bibitem{BPR}
S.~Basu, R.~Pollack, and M.-F.~Roy,
\emph{Algorithms in Real Algebraic Geometry}, second edition,
Algorithms and Computation in Mathematics \textbf{10}, Springer, 2006.
\href{https://doi.org/10.1007/3-540-33099-2}{doi:10.1007/3-540-33099-2}.

\bibitem{BFZ}
P.~V.~M. Blagojevi\'c, F.~Frick, and G.~M. Ziegler,
\emph{Tverberg plus constraints},
Bull. Lond. Math. Soc. \textbf{46} (2014), 953--967.
\href{https://doi.org/10.1112/blms/bdu049}{doi:10.1112/blms/bdu049}.

\bibitem{Frick}
F.~Frick,
\emph{Counterexamples to the topological Tverberg conjecture},
Oberwolfach Rep. \textbf{12} (2015), 318--321;
\href{https://arxiv.org/abs/1502.00947}{arXiv:1502.00947}.

\bibitem{Gromov}
M.~Gromov,
\emph{Singularities, expanders and topology of maps. Part 2:
From combinatorics to topology via algebraic isoperimetry},
Geom. Funct. Anal. \textbf{20} (2010), 416--526.
\href{https://doi.org/10.1007/s00039-010-0073-8}{doi:10.1007/s00039-010-0073-8}.

\bibitem{Hatcher}
A.~Hatcher, \emph{Algebraic Topology}, Cambridge University Press, 2002.

\bibitem{MW}
I.~Mabillard and U.~Wagner,
\emph{Eliminating higher-multiplicity intersections, I.
A Whitney trick for Tverberg-type problems},
\href{https://arxiv.org/abs/1508.02349}{arXiv:1508.02349} (2015).

\bibitem{LZ}
E.~Le\'on and G.~M.~Ziegler,
\emph{Spaces of convex $n$-partitions}, arXiv:1511.02904 (2015).
\end{thebibliography}
\end{document}
